\documentclass[12pt
,a4paper
,titlepage
,twoside
%,openany % ouvre les sections indifféremment sur la page de gauche ou de droite
]{book}

% Paramétrages de la langue et de l'encodage
\usepackage[utf8]{inputenc}
\usepackage[french]{babel}
\usepackage[T1]{fontenc}

% formules mathématiques
\usepackage{amsmath}
\usepackage{amsfonts}
\usepackage{amssymb}

% gestion des graphiques
\usepackage{graphicx}

% Gestion des couleurs
\usepackage[usenames,dvipsnames]{xcolor}

% Redefinition des liens web
\usepackage{url}
\usepackage[colorlinks=false,urlbordercolor=white,linkbordercolor=white]{hyperref}

% Définition des marges de la page
\usepackage[left=2cm,right=2cm,top=2cm,bottom=2cm]{geometry}

% Polices de caractères
\usepackage{lmodern}
\usepackage{mathptmx} % times, y compris dans les formules mathématiques

% Insertion de code source dans le texte, à utiliser avec \begin{lstlisting} et \lstset{java|html|php...}
\usepackage{listings}

% Définition des entêtes
\usepackage{fancyhdr}
\pagestyle{fancy}

% Redéfinition des titres de section
\usepackage{titlesec}

% Insertion de graphiques à des emplacements définis
\usepackage[abs]{overpic}

% règles typographiques de l'Imprimerie nationale
\usepackage[all]{nowidow}
\usepackage[
% frenchchapters renomme le premier chapitre, mais :
%	- cela pose problème dans la table des matières
%	- cela ne peut être utilisé qu'avec la renumérotation des chapitres activée
%frenchchapters,
parindent,
lastparline,
hyphenation
]{impnattypo}

% Renumérotation des chapitres
%\renewcommand{\thesection}{\Alph{section})}
%\renewcommand{\thesubsection}{\arabic{subsection} -}
%\renewcommand{\thesubsubsection}{\alph{subsubsection} -}
%\usepackage{engrec}
%\renewcommand{\theparagraph}{\engrec{paragraph})}
%\setcounter{secnumdepth}{4}

% Génération du code Ipsum lorem
\usepackage{blindtext}



% Definition des couleurs IRSTEA
\definecolor{titreColor}{RGB}{0,58,128}  % Marine
\definecolor{stitreColor}{RGB}{0,158,224}  % Ocean
\definecolor{auteurColor}{RGB}{0,58,128}     % Marine
\definecolor{texteColor}{RGB}{164,196,0}     % Prairie

% Definition des chapitres
\titleformat{\chapter}[display]
{\normalfont\Large\filcenter\sffamily}
%{\titlerule[1pt]%
% \vspace{1pt}%
% \titlerule
% \vspace{1pc}%
{
 \Large\color{titreColor}{
 %\MakeUppercase{\chaptertitlename}
 \chaptertitlename~\thechapter}
 }
{1pc}
%{\titlerule
%% \vspace{1pc}%
 \Huge

\titleformat{\section}
{\color{titreColor}\normalfont\Large\bfseries\sffamily}
{\color{titreColor}\thesection}{1em}{}

\titleformat{\subsection}
{\color{stitreColor}\bfseries\sffamily}
{\color{stitreColor}\thesubsection}{1em}{}

% Bibliographie
% natbib est indispensable si la biblio contient des accents
% Options pour natbib (extrait de http://merkel.zoneo.net/Latex/natbib.php)
%    round: (par défaut) pour des parenthèses arondies (());
%    square: pour des crochets ([]);
%    curly: pour des accolades ({});
%    angle: pour des équerres (<>) ;
%    colon: (par défaut) pour séparer les citations multiples par deux points (:);
%    comma: pour utiliser une virgule comme séparateur;
%    authoryear: (par défaut) pour des citations auteurs-année;
%    numbers: pour des citations numériques;
%    super: pour des citations numériques en exposant, comme dans Nature;
%    sort: ordonne les citations multiples dans l'ordre dans lequel elles apparaissent dans la bibliographie;
%    sort&compress: comme sort mais en plus les citations numériques multiples sont comprimées, si possible (3-6, 15, par exemple);
%    longnamesfirst: transforme la première citation à une référence en une version étoilée (avec la liste complète des auteurs) et le citations suivantes normales (liste abbrégée);
%    sectionbib: pour redéfinir \thebibliography pour avoir une \section* à la place d'un \chapter*; valide seulement pour les classes de document possédant la commande \chapter; à utiliser avec le paquetage  chapterbib;
%    nonamebreak: garde tous les noms d'auteurs d'une citation sur une même ligne; celà cause des problèmes de débordement, mais permet de résoudre certains problèmes liés à hyperref.

% En cas de souci, supprimez les fichiers .aux et .bbi après modification des paramètres
\usepackage[square,sort,comma,numbers]{natbib}
% styles natbib natifs : abbrvnat, plainnat, unsrtnat
\bibliographystyle{unsrtnat}
\usepackage{hypernat}
% Ajout de la référence à la bibliographie dans la table des matières
\usepackage[nottoc, notlof, notlot]{tocbibind}

%Données de titre et d'auteur pour la page de garde
\newcommand{\titre}{Logiciel COLLEC}
\newcommand{\sousTitre}{Installation et configuration}
\newcommand{\auteur}{Éric Quinton}
\newcommand{\dateModif}{\today}


\begin{document}
%Supprime les veuves et orphelines
\widowpenalty=10000
\clubpenalty=10000
\raggedbottom 

% Integrer la page de garde
%\setcounter{page}{0}
\thispagestyle{empty}
% Logo IRSTEA
\vspace{-2cm}
\hspace{-2cm}
\includegraphics[width=3.06cm,height=3.06cm,keepaspectratio]{logo_irstea}%


\vspace*{4cm}
\hspace{5cm}
\setlength\unitlength{1mm}
% Logo de titre
\begin{overpic}[width=9.44cm,height=9.57cm,keepaspectratio]{logo_fond_droite.png}
% Titre du document
\put(-50,50){
\begin{minipage}{0.7\linewidth}
\Huge\flushright \color{titreColor}{\bfseries\sffamily\titre{}}\\
% sous-titre du document
\color{stitreColor}{\Large \bfseries\sffamily\sousTitre{}}
\end{minipage}
}
\end{overpic}

%  date et auteur
% creation de l'espace a gauche
\vspace*{1cm}
\begin{minipage}{0.5\linewidth}
\hfill
\end{minipage}
% positionnement
\begin{minipage}{0.5\linewidth}\flushleft{
% Date
\textcolor{auteurColor}{\Large\sffamily\dateModif{}}\\
\vspace*{0.1cm}
% Auteur
\textcolor{auteurColor}{\Large\sffamily\auteur{}}\\
\vspace{0.5cm}
% Adresse
\textcolor{texteColor}{\sffamily\textbf{IRSTEA} - Centre de Bordeaux\\
50, avenue de Verdun, Gazinet\\
33612 CESTAS Cedex }
}
\end{minipage}
% Ligne de logos
\begin{minipage}{\linewidth}

% logos complémentaires
%\vspace{3cm}
%\hspace{2cm}
%\includegraphics[width=3cm,height=3cm,keepaspectratio]
%{emplacement_logo}%
%\vspace{-1cm}
%\hspace{1cm}
%\includegraphics[width=3cm,height=3cm,keepaspectratio]
%{emplacement_logo}%
%\vspace{-1cm}
%\hspace{1cm}
%\includegraphics[width=3cm,height=3cm,keepaspectratio]
%{emplacement_logo}%
\end{minipage}
% Définition des entêtes
\fancyhead{}
\fancyhead[CO]{\leftmark\sffamily}
\fancyhead[CE]{ \sffamily\titre{}}
\fancyfoot[CO]{\sffamily\thepage}
\fancyfoot[CE]{\sffamily\thepage}
% Redéfinition de \cleardoublepage pour créer une page totalement vide
\makeatletter
\def\cleardoublepage{\clearpage\if@twoside \ifodd\c@page\else
  \hbox{}
  \vspace*{\fill}

  \vspace{\fill}
  \thispagestyle{empty}
  \newpage
  \if@twocolumn\hbox{}\newpage\fi\fi\fi}
\makeatother

% \cleardoublepage permet de générer une page vide 
% si le chapitre ne commence pas sur la page de droite

% Ajout d'un préambule
\frontmatter
\cleardoublepage
\chapter{Préambule}
\blindtext

% Début réel du texte
\mainmatter
%\cleardoublepage
\chapter{Introduction}
\section{Première section}
Référence bibliographique au livre \citet{livre}, à l'article \citet{article}, à la thèse \citep{these}.

La bibliographie utilise le fichier \textit{irstea.bib} dans cet exemple.
\subsection{sous-section}
\blindtext
\subsection{autre sous-section}
\blindtext

\blindtext

\blindtext

% Second chapitre
%\cleardoublepage
\chapter{Second chapitre}
% \input{chapitre2}

%Annexes
\appendix
\chapter{Première annexe}

%Bibliographie
\backmatter
% Integration de la biblio
% Pour insérer toutes les références : 
%\nocite{*}
% Pour intégrer une référence non citée : 
%\nocite{ref}
\nocite{*}
\bibliography{collec}

% Table des matières
\tableofcontents
\end{document}